%\usepackage[utf8]{inputenc}
%\usepackage[T1]{fontenc}
\usepackage[spanish,es-tabla,es-noshorthands]{babel}
\usepackage{csquotes}                 


% GEOMETRÍA Y FORMATO
\usepackage[
    a4paper,
    left=2.5cm,
    right=2.5cm,
    top=3cm,
    bottom=3cm
]{geometry}

\setlength{\parindent}{0pt}
\setlength{\parskip}{0.5\baselineskip plus 2pt minus 1pt}     
\usepackage{setspace}
\setstretch{1.25}                     

\renewcommand{\thefootnote}{\roman{footnote}}  

% COLORES
\usepackage[dvipsnames,table,xcdraw]{xcolor}

% GRÁFICOS Y TABLAS
\usepackage{graphicx}
\usepackage{float}
\usepackage{caption}
\usepackage{subcaption}
\usepackage{booktabs}                 
\usepackage{array}
\usepackage{multirow}
\usepackage{tabularx}
\usepackage{xltabular}
\usepackage{makecell}
\usepackage{wrapfig} 
\usepackage{verbatim} 
\usepackage{comment}    
\usepackage{pgfgantt}    

\usepackage{array}

\newcolumntype{L}[1]{>{\raggedright\arraybackslash}p{#1}}
\newcolumntype{C}[1]{>{\centering\arraybackslash}p{#1}}
\newcolumntype{R}[1]{>{\raggedleft\arraybackslash}p{#1}}

% TIKZ Y DIAGRAMAS
\usepackage{tikz}                     % declaracion explicita (best practice)
\usepackage{pgfplots}
\pgfplotsset{compat=1.18}
\usetikzlibrary{
    positioning,                      % right=of, below=of, etc.
    arrows.meta,                      % puntas modernas (Stealth, Latex, etc.)
    shapes.geometric,                 % diamond, ellipse, trapezoid
    shapes.misc,                      % chamfered rectangle, etc.
    calc,                             % aritmetica de coordenadas ($(a)+(1,0)$)
    fit,                              % agrupar nodos en una caja
    backgrounds,                      % capas de fondo (\begin{scope}[on background layer])
    decorations.pathreplacing         % llaves y anotaciones sobre flechas
}


% Estilos comunes para diagramas (definidos UNA vez aqui, usados en todas las figuras)
\tikzset{
    block/.style={
        rectangle, draw, rounded corners=2pt,
        minimum height=9mm, minimum width=24mm,
        align=center, font=\small,
        inner sep=4pt
    },
    sensor/.style={
        block, fill=background_light
    },
    decision/.style={
        diamond, draw, aspect=2,
        align=center, font=\small,
        inner sep=1pt
    },
    terminal/.style={
        ellipse, draw, minimum width=20mm, minimum height=8mm,
        align=center, font=\small
    },
    arrow/.style={
        -Stealth, thick
    },
    bus/.style={
        -Stealth, very thick, double
    },
    label/.style={
        font=\scriptsize, midway, fill=white, inner sep=1pt
    }
}

% Colores personalizados
\definecolor{background_light}{RGB}{245,245,245}
\definecolor{border_color}{RGB}{120,120,120}
\definecolor{keyword_blue}{RGB}{0,0,180}
\definecolor{comment_green}{RGB}{34,139,34}
\definecolor{string_red}{RGB}{163,21,21}
\definecolor{number_gray}{RGB}{128,128,128}


% MATEMÁTICAS Y UNIDADES
\usepackage{amsmath,amsfonts,amssymb,amsthm}
\usepackage{textcomp}
\usepackage{gensymb}                 
\usepackage{siunitx}                  
\sisetup{
    output-decimal-marker={,},        
    per-mode=symbol,
    range-phrase={\,--\,},
    range-units=single,
    detect-all
}
% Definir unidades adicionales
\DeclareSIUnit\baud{Bd}
\DeclareSIUnit\sample{S}
\DeclareSIUnit{\eur}{\text{\euro}}

% LISTADOS DE CÓDIGO
\usepackage{listings}
\usepackage[framemethod=tikz]{mdframed}

% Marco para bloques de código
\mdfdefinestyle{CodeFrame}{
    roundcorner=5pt,
    backgroundcolor=background_light,
    linecolor=border_color,
    linewidth=0.5pt,
    innertopmargin=10pt,
    innerbottommargin=10pt,
    innerleftmargin=10pt,
    innerrightmargin=10pt,
    skipabove=10pt,
    skipbelow=10pt,
}

% Configuración base para código
\lstset{
    basicstyle=\ttfamily\small,
    breaklines=true,
    breakatwhitespace=false,
    columns=flexible,
    keepspaces=true,
    showstringspaces=false,
    showspaces=false,
    tabsize=4,
    frame=none,
    numbers=left,
    numberstyle=\tiny\color{number_gray},
    numbersep=8pt,
    xleftmargin=15pt,
    captionpos=b,
    extendedchars=true,
    inputencoding=utf8,
    literate=
        {á}{{\'a}}1 {é}{{\'e}}1 {í}{{\'i}}1 {ó}{{\'o}}1 {ú}{{\'u}}1
        {Á}{{\'A}}1 {É}{{\'E}}1 {Í}{{\'I}}1 {Ó}{{\'O}}1 {Ú}{{\'U}}1
        {ñ}{{\~n}}1 {Ñ}{{\~N}}1 {ü}{{\"u}}1 {Ü}{{\"U}}1,
}

% Estilo Arduino/C++ (para firmware ESP32)
\lstdefinestyle{Arduino}{
    language=C++,
    keywordstyle=\color{keyword_blue}\bfseries,
    commentstyle=\color{comment_green}\itshape,
    stringstyle=\color{string_red},
    morekeywords={
        HIGH,LOW,INPUT,OUTPUT,INPUT_PULLUP,LED_BUILTIN,
        Serial,Wire,String,boolean,byte,word,
        pinMode,digitalWrite,digitalRead,analogRead,analogWrite,
        delay,millis,micros,setup,loop,
        Serial1,Serial2,HardwareSerial,BluetoothSerial,
        begin,print,println,available,read,write,flush,
        uint8_t,uint16_t,uint32_t,int8_t,int16_t,int32_t,
        float,double,void,bool,true,false,NULL,
        PROGMEM,F,sizeof,static_cast
    },
    morecomment=[l]{//},
    morecomment=[s]{/*}{*/},
    morestring=[b]",
    morestring=[b]',
}

% Estilo Processing/Java (para software PFD)
\lstdefinestyle{Processing}{
    language=Java,
    keywordstyle=\color{keyword_blue}\bfseries,
    commentstyle=\color{comment_green}\itshape,
    stringstyle=\color{string_red},
    morekeywords={
        void,setup,draw,settings,size,fullScreen,
        background,fill,stroke,noFill,noStroke,strokeWeight,
        ellipse,rect,line,arc,triangle,quad,point,vertex,
        beginShape,endShape,curveVertex,bezierVertex,
        translate,rotate,scale,pushMatrix,popMatrix,resetMatrix,
        color,loadImage,image,PImage,PFont,textFont,createFont,
        text,textSize,textAlign,textWidth,
        mouseX,mouseY,mousePressed,keyPressed,key,keyCode,
        width,height,frameRate,frameCount,displayWidth,displayHeight,
        Serial,printArray,println,print,nf,nfc,nfp,nfs,
        map,constrain,lerp,norm,dist,mag,
        radians,degrees,sin,cos,tan,atan2,asin,acos,
        sqrt,pow,abs,min,max,floor,ceil,round,
        random,noise,randomSeed,noiseSeed,
        CENTER,LEFT,RIGHT,TOP,BOTTOM,CORNER,CORNERS,BASELINE,
        PI,TWO_PI,HALF_PI,QUARTER_PI,TAU,
        RGB,HSB,ARGB,ALPHA,
        CLOSE,OPEN,CHORD,PIE
    },
    morecomment=[l]{//},
    morecomment=[s]{/*}{*/},
}

% Estilo Python
\lstdefinestyle{Python}{
    language=Python,
    keywordstyle=\color{keyword_blue}\bfseries,
    commentstyle=\color{comment_green}\itshape,
    stringstyle=\color{string_red},
    morekeywords={self,True,False,None,as,with,lambda,yield,async,await},
}

% Nombre de los listados en español
\renewcommand{\lstlistingname}{Código}
\renewcommand{\lstlistlistingname}{Índice de Códigos}


% LISTAS Y ENUMERACIÓN
\usepackage{enumitem}


% ENCABEZADOS Y PIES DE PÁGINA
\usepackage{fancyhdr}
\setlength{\headheight}{15pt}
\setlength{\headsep}{0.5cm}

\pagestyle{fancy}
\fancyhf{} 

% Estilo para páginas normales
\fancyhead[L]{\small\leftmark}                    
\fancyhead[R]{\thepage} 
\renewcommand{\headrulewidth}{0.4pt}
\renewcommand{\footrulewidth}{0pt}

% Estilo para páginas de inicio de capítulo
\fancypagestyle{plain}{
    \fancyhf{}
    \fancyfoot[C]{\thepage}
    \renewcommand{\headrulewidth}{0pt}
}


% BIBLIOGRAFÍA
% BIBLIOGRAFIA (UNE-ISO 690, sistema numerico)
\usepackage[
    backend=biber,
    style=iso-numeric,      % <-- biblatex-iso690: implementa UNE-ISO 690
    sorting=none,           % orden de aparicion en el texto
    maxnames=3,
    minnames=1,
    shortnumeration=true,   % edicion abreviada: "3.a ed." en vez de "tercera edicion"
    pagetotal=false,        % no mostrar total de paginas
]{biblatex}

\addbibresource{../Common/references.bib}


% HIPERVÍNCULOS 
\usepackage{xurl}                     
\usepackage[
    colorlinks=true,
    linkcolor=black,
    urlcolor=blue,
    citecolor=blue,
    breaklinks=true,
    pdfauthor={Arnau},
    pdftitle={Flight Instruments Emulation Platform},
    pdfsubject={TFG - Trabajo de Fin de Grado},
    pdfkeywords={ESP32, IMU, GPS, PFD, instrumentación aérea},
    bookmarks=true,
    bookmarksnumbered=true
]{hyperref}


% REFERENCIAS CRUZADAS
\usepackage[spanish,capitalise]{cleveref}

% Personalizar nombres
\crefname{figure}{Figura}{Figuras}
\crefname{table}{Tabla}{Tablas}
\crefname{equation}{Ecuación}{Ecuaciones}
\crefname{chapter}{Capítulo}{Capítulos}
\crefname{section}{Sección}{Secciones}
\crefname{appendix}{Apéndice}{Apéndices}
\crefname{listing}{Código}{Códigos}


% SÍMBOLOS Y OTROS
\usepackage{eurosym}
\usepackage{pdfpages}                 


% COMANDOS PERSONALIZADOS

% Código inline: \code{setup()}
\newcommand{\code}[1]{\texttt{#1}}

% Nombres de archivo: \file{main.ino}
\newcommand{\file}[1]{\texttt{#1}}

% Componentes: \component{ESP32}
\newcommand{\component}[1]{\textsc{#1}}

% Registros: \register{PWR\_MGMT\_1}
\newcommand{\register}[1]{\texttt{#1}}

% Dirección I2C: \iicaddr{0x68}
\newcommand{\iicaddr}[1]{\texttt{#1}}

% Pin: \pin{GPIO21}
\newcommand{\pin}[1]{\texttt{#1}}

% Abreviatura de pendiente
\newcommand{\pte}{pte.}